% Cardiovascular Research: neutral adaptation of an inspected 2023 article.
% Original prose, fictional authors and native diagrams; no source-paper assets.
\PassOptionsToPackage{no-math}{fontspec}
\documentclass[UTF8,fontset=none,scheme=plain]{ctexart}
\usepackage[paperwidth=595.276bp,paperheight=793.701bp,left=25.285bp,right=29.991bp,top=28bp,bottom=42bp,headheight=22bp,headsep=16.255bp,footskip=22bp]{geometry}
\newcommand{\JournalName}{Cardiovascular Research}
% Shared dependencies and neutral material. Compile with XeLaTeX.
% No publisher logos, paper text, data, figures, or download metadata are included.
\usepackage{amsmath,amssymb,graphicx,xcolor,tikz,fancyhdr,needspace,xurl,hyperref,balance,etoolbox}
% The legacy balance package omits the first split box's final baseline depth.
% Reserve that measured depth rather than hiding a CJK overfull-box warning.
\makeatletter
\patchcmd{\@BAlancecol}{\setbox\@leftcolumn=\vbox to \@tempdima}{\advance\@tempdima by \dp1\relax\setbox\@leftcolumn=\vbox to \@tempdima}{}{\PackageWarning{cns-template}{Balance depth patch unavailable; inspect final columns}}
\makeatother
\usetikzlibrary{arrows.meta,positioning}
\usepackage[protrusion=true,expansion=false]{microtype}
\IfFontExistsTF{Noto Serif CJK SC}{%
  \setCJKmainfont{Noto Serif CJK SC}%
}{%
  \setCJKmainfont[BoldFont=FandolSong-Bold.otf,ItalicFont=FandolKai-Regular.otf]{FandolSong-Regular.otf}%
}
\IfFontExistsTF{Noto Sans CJK SC}{%
  \setCJKsansfont{Noto Sans CJK SC}%
}{%
  \setCJKsansfont[BoldFont=FandolHei-Bold.otf]{FandolHei-Regular.otf}%
}
\setmainfont{TeX Gyre Termes}\setsansfont{TeX Gyre Heros}
\setmonofont{TeX Gyre Cursor}
\definecolor{JournalBlue}{HTML}{0097C9}
\definecolor{JournalRed}{HTML}{B83443}
\definecolor{ScienceRed}{HTML}{CD1021}
\definecolor{CiteBlue}{HTML}{0000EE}
\hypersetup{colorlinks=true,linkcolor=JournalBlue,urlcolor=JournalBlue,
  pdftitle={CNS论文汉化：中性版式示例},pdfauthor={示例模板；非期刊官方模板},
  pdfsubject={Neutral source-inspired Chinese layout demonstration; no real research claims}}
\urlstyle{same}\raggedbottom
\setlength{\parindent}{1em}\setlength{\parskip}{0pt}
\setlength{\emergencystretch}{2em}\tolerance=2200
\XeTeXlinebreaklocale "zh"\XeTeXlinebreakskip=0pt plus .1pt
\pagestyle{fancy}\fancyhf{}
\renewcommand{\headrulewidth}{0pt}\renewcommand{\footrulewidth}{0pt}
\newcommand{\PaperTitle}{中文题名占位：从组织结构到证据解释}
\newcommand{\SampleAuthors}{Author A\textsuperscript{1}, Author B\textsuperscript{2,*}}
\newcommand{\SampleAffiliations}{\textsuperscript{1}作者单位一（示例占位）。\quad
  \textsuperscript{2}作者单位二（示例占位）。\par
  *通信作者信息待填。DOI、收稿和出版日期均应从目标原文核实。}
\newcommand{\SampleAbstract}{本段是中性摘要占位文字，不陈述任何真实研究发现。它演示中文题名、作者信息、摘要和双栏正文之间的层级关系。实际使用时，应完整替换为目标论文的忠实中文译文，并保留研究对象、方法、结果和结论的限定条件。图注、参考文献、作者声明和来源信息均属于独立的覆盖项。模板中的图形仅说明排版结构，不表示生物学机制、实验结果或统计关系。}
\newcommand{\DemoCaption}{\textbf{图1. 中性证据链示意（占位图，无实验数据）。}
  \textbf{A}，输入节点表示需要核对的原文材料。
  \textbf{B}，处理节点表示逐段转换及复核。
  \textbf{C}，输出节点表示正文、图注和书目等独立交付部分。
  箭头仅表示本示例的阅读顺序，不代表因果关系、效应方向或数量变化。
  该图由模板自行绘制，所有节点均为文字占位；无样本量、误差线、统计检验或比例尺。
  实际论文图注应逐项保留原稿报告的样本量、单位、检验方法、限定及图形署名。}
\newcommand{\DemoVector}{%
 \resizebox{\linewidth}{!}{%
 \begin{tikzpicture}[font=\sffamily\small,node distance=18mm,
   box/.style={draw=black!70,rounded corners=2pt,fill=black!3,minimum width=26mm,minimum height=16mm,align=center},
   flow/.style={-{Stealth[length=2.5mm]},line width=.7pt}]
  \node[box] (a) {\textbf{A}\quad 输入\\原文材料};
  \node[box,right=of a] (b) {\textbf{B}\quad 处理\\转换与复核};
  \node[box,right=of b] (c) {\textbf{C}\quad 输出\\中文译稿};
  \draw[flow] (a)--(b);\draw[flow] (b)--(c);
  \node[below=11mm of a,align=center] {标题与摘要\\正文与图表};
  \node[below=11mm of b,align=center] {保留限定\\检查单位};
  \node[below=11mm of c,align=center] {可追溯引文\\明确原文信息};
  \node[below=29mm of b,font=\sffamily\footnotesize,text=black!65] {示意占位；不包含真实实验数据};
 \end{tikzpicture}}%
}
\newcommand{\SampleBibliography}{%
 \section*{参考文献（示例占位）}
 {\fontsize{7.6bp}{10bp}\selectfont\setlength{\parindent}{0pt}\interlinepenalty=10000
 \hypertarget{ref:1}{}1. Author A. (20XX). Example reference title. \textit{Journal placeholder}. 卷、页码与DOI待填。此条仅供演示引文跳转。\par\smallskip
 \hypertarget{ref:2}{}2. Author B. (20YY). Another example reference title. \textit{Journal placeholder}. 出版信息待填；不要把占位条目带入正式译稿。\par}
}
\newcommand{\SampleSourceInfo}{%
 \section*{原文信息}
 以下字段全部为示例占位：作者；英文原标题；期刊；年份；卷期；页码或文章号；DOI或来源链接。实际使用时必须逐项替换并核对。\par
 本文件是\JournalName 风格的非官方中文译稿样板，依据一次已核验文章版式进行结构化改编，不代表该期刊所有年代或文体的统一标准，也不保证像素级复刻。无真实论文内容或期刊标识图像。\par
}
% Deliberately neutral demo prose, not translated from a published paper.
\newcommand{\SampleOpening}{%
 \section*{引言：示例正文与引文}
 这是一段用于检查双栏阅读节奏的中性示例文字。段落中的术语、缩写和标点需要在实际项目中与源文保持一致。中文译稿应呈现原论文的完整论述，而不是把多段证据压缩成一段摘要。这里的引文只链接到本文件末尾明确标为占位的参考文献\SampleCiteOne。\par
 若原句仅表示相关、可能或提示，译文也应保留相同的证据强度。数字、单位、正负号、上下标和范围端点应逐项核对。原文存在内部矛盾时，应分别保留相应位置的陈述，再用简短译注说明；不能根据常识把不同的数字静默统一。\par
 \subsection*{内容层级与跨页衔接}
 此小标题演示二级层级。实际论文可能采用不同的小标题密度，模板不会替代对原文结构的识别。双栏文章需按左栏、右栏和跨页续句的顺序读取。若段落在页末中断，应在完成整句核对后接入下一页，避免把图中文字误当正文，或重复翻译同一段。\par
 图表具有独立的信息范围。正文引用到某幅图，并不意味着其图注可以省略。图注中的每个分面、比较对象、时间点、样本量及统计说明都应对应源文。英文图内标签可以保留，并在图旁或文末提供中文释义；实际采用哪种安排，需要依据图形尺寸和译文长度判断\SampleCiteTwo。\par
 \subsection*{来源与可维护数据}
 本段继续演示正文密度。可编辑文本和图注数据应与最终版面分开维护，以便修订时准确定位变化。若只调整版式，应能证明正文、图注和书目内容没有被意外改写。更换字体或纸张后，须重新检查断行、图注与图片的相对位置以及页脚可读性。\par
 文内引用与书目之间应有明确对应关系。这里的两条演示书目没有真实DOI，也没有虚构的实验信息。使用模板时可增删条目，但必须同时更新编号、链接目标和引用位置。页码只表示中文译稿自己的分页，不用于冒充原出版页码。\par
 \subsection*{结果叙述的演示占位}
 本节的文字仅用于呈现长段落在双栏页面中的流动，不报告真实结果。实际翻译应保留观察对象、实验条件和结局指标之间的关系。例如，原文分别描述体外实验与动物实验时，不能把两类结果合并为同一层面的结论；研究使用的对照、时间顺序和分析范围也必须保留。图表中的比较范围不应扩大为对整个人群或全部疾病阶段的判断。\par
 论述中的转折与限制同样属于正文。若研究先提出一种解释，再讨论另一种可能，译文应沿用这一推进顺序。作者明确写出的不确定性应能在中文中被识别；如果尚未确定某个现象的方向或机制，就不能改成已经建立的因果关系。对于术语近似但研究层级不同的概念，可在首次出现时保留缩写，后文保持统一。\par
 \subsection*{方法信息与可复核性}
 方法部分往往包含密集数字和专有名词。翻译时须区分浓度、总量、频率和处理时长，不能因为中文句式调整而改变参数的归属。试剂名称、软件版本和数据访问标识通常保留原样；上下限、排除规则和显著性标准则需要逐字符核对。这里没有提供任何真实实验参数，避免把模板占位误认为可以直接执行的方法。\par
 表格与补充材料的引用也需要保留原层级。正文对补充图的指向并不意味着模板中已经附有该补充图；若源文件只包含主文，应如实标注所交付的范围。图表编号、附录顺序和参考文献数量都应由目标文章决定，本示例中的结构不构成固定配额，也不要求所有论文采用同样的章节。\par
}
\newcommand{\SampleClosing}{%
 \section*{讨论：版式适配与范围}
 上方矢量图是模板原创的排版占位图，不来自任何论文。图形与完整演示图注被放在同一页，便于读者直接核对分面说明。真实项目应使用经过授权或用户提供的原图素材，并保留图注署名与许可信息。尺寸不足时，可适度缩放或调整栏宽，不以删除图注解决空间问题。\par
 模板仅记录一次已核验文章的可观察设计特征。即使文章出自同一期刊，不同年代、研究论文与综述、校正清样与正式版本也可能采用不同的页眉和引文形式。应以用户提供的实际PDF为准，不能仅凭期刊名称推断配色、字号或引用方式\SampleCiteOne。\par
 \subsection*{演示检查项}
 本样板用于展示可编译性、中文字形、标题层级、跨栏图注、页眉页脚和内部链接。图中没有实际数据，因此不应被当作科学插图引用。替换内容后，还需检查所有页的渲染结果，确认无截字、碰撞、异常空白和遗漏段落。\par
 \SampleBibliography
 \SampleSourceInfo
}


\usepackage{multicol,tabularx,array,colortbl}
\definecolor{CVRRed}{rgb}{.935,.28,.21}
\definecolor{CVRCream}{rgb}{.995,.925,.895}
\definecolor{CVRGrey}{rgb}{.9,.905,.905}
\definecolor{CVRCiteBlue}{HTML}{0000FF}
\setlength{\columnsep}{15.488bp}\raggedcolumns
\renewcommand{\familydefault}{\sfdefault}
\renewcommand{\normalsize}{\fontsize{9.1bp}{12.1bp}\selectfont}
\setlength{\parindent}{1em}\setlength{\parskip}{2bp}
\ctexset{section={format=\sffamily\bfseries\color{CVRRed}\fontsize{15bp}{18bp}\selectfont,beforeskip=11bp,afterskip=5bp},
 subsection={format=\sffamily\bfseries\color{black}\fontsize{12.5bp}{15bp}\selectfont,beforeskip=8bp,afterskip=4bp}}
\hypersetup{colorlinks=true,linkcolor=CVRCiteBlue,urlcolor=CVRCiteBlue,
 pdftitle={Cardiovascular Research风格：原创中性中文排版示例},pdfauthor={原创中性模板；非期刊官方模板},
 pdfsubject={Original neutral example; no source-paper prose, figures or data}}
\fancyhead[L]{\sffamily\fontsize{8.5bp}{10bp}\selectfont 资料组织与记录对应（原创中性演示）}
\fancyhead[R]{\sffamily\bfseries\fontsize{10bp}{12bp}\selectfont\thepage}
\renewcommand{\headrulewidth}{.5bp}
\renewcommand{\headrule}{\hbox to\headwidth{\color{CVRRed}\leaders\hrule height\headrulewidth\hfill}}
\fancyfoot[L]{\sffamily\fontsize{6.8bp}{9bp}\selectfont Cardiovascular Research风格｜非官方中文样板}
\fancyfoot[R]{\sffamily\fontsize{6.8bp}{9bp}\selectfont 中文译页\thepage｜仅供学术交流}
\newcommand{\CVRCite}[1]{\textsuperscript{\hyperlink{cvr-ref:#1}{#1}}}
\newcommand{\CVRDotRule}{\par\vspace{3bp}\noindent\textcolor{CVRRed}{\leaders\hbox{\kern1bp.\kern1bp}\hfill\kern0pt}\par\vspace{4bp}}
\newcommand{\CVRRow}[2]{\noindent\begin{minipage}[t]{76bp}\vspace{0pt}\bfseries\color{CVRRed}\raggedright #1\end{minipage}\hspace{9bp}\begin{minipage}[t]{\dimexpr\linewidth-85bp\relax}\vspace{0pt}#2\end{minipage}\par}
\newsavebox{\CVRSummaryBox}
\newenvironment{cvrsummary}{\par\begingroup\setlength{\fboxsep}{0pt}\begin{lrbox}{\CVRSummaryBox}\begin{minipage}{\textwidth}\vspace{7bp}\fontsize{9.1bp}{12.1bp}\selectfont}{\vspace{8bp}\end{minipage}\end{lrbox}\noindent\makebox[\textwidth][l]{\hspace*{-25.285bp}\colorbox{CVRCream}{\hspace*{25.285bp}\usebox{\CVRSummaryBox}\hspace*{29.991bp}}}\endgroup\par}
\newcommand{\CVRBody}{本段为原创中性排版文字，不报告任何真实研究。资料整理应区分正文、方法、图表、书目和声明，先恢复跨栏与跨页的阅读顺序，再核对记录之间的对应关系。中文自然重排可以改变页数，不能改变比较对象、否定、单位或证据范围。图表是独立内容单元，完整图注及全部表脚注均须保留。}
\begin{document}\sffamily\normalsize
\fancypagestyle{cvrfirst}{\fancyhead{}\renewcommand{\headrulewidth}{0pt}}
\thispagestyle{cvrfirst}
\noindent{\bfseries\fontsize{17bp}{21bp}\selectfont Cardiovascular Research\par}
{\fontsize{8bp}{10bp}\selectfont 原刊年份、卷期、页码、DOI及版本信息按目标来源填写\par}\vspace{10bp}
{\color{CVRRed}\hrule height .5bp}\vspace{17bp}
\noindent{\bfseries\fontsize{23bp}{28bp}\selectfont 从资料组织到证据对应：\\不同记录的整合与可追溯呈现\par}\vspace{12bp}
\noindent{\bfseries\fontsize{12bp}{15bp}\selectfont Author Placeholder A\textsuperscript{1,*}, Author Placeholder B\textsuperscript{2}\par}\vspace{6bp}
{\fontsize{7.5bp}{10bp}\selectfont\noindent\textsuperscript{1}原创中性作者单位一；\textsuperscript{2}原创中性作者单位二。通信信息按所给原文填写，此处没有真实作者、地址或邮箱。\par}
\vspace{7bp}{\fontsize{8bp}{10bp}\selectfont 收稿、修回、接受及提前在线日期：按当前来源逐项核验。\par
\textbf{首次审稿用时：}按原文填写，不等同于出版日期。\par}\vspace{13bp}
\begin{cvrsummary}
\CVRRow{Aims\newline 目的}{本段演示目的摘要，不报告真实研究。完整记录的整理需要明确来源范围，区分材料清单与解释性文字，并保留各项记录的上下文。本示例仅说明排版，不设置研究对象、样本量或科学结论。}
\CVRDotRule
\CVRRow{Methods and\newline results\newline 方法与结果}{本段演示合并的“方法与结果”标签及其正文。实际稿应保留原作者采用的方法、比较对象、结果方向、限定条件和不确定性；不把这一合并字段强改成其他期刊的结构。正文、图表、方法和声明分别建立对应关系。原文提供的参数、单位、统计信息和数据来源均应逐项核对，本例不包含真实实验流程或定量结果。}
\CVRDotRule
\CVRRow{Conclusions\newline 结论}{本段是结论结构占位。正式汉化应保持原文证据强度及适用范围；中文分页可以变化，结论尾句和限制条件仍须完整保留。}
\CVRDotRule
\CVRRow{Keywords\newline 关键词}{中性示例\enspace\textbullet\enspace 资料组织\enspace\textbullet\enspace 记录对应}
\end{cvrsummary}\vspace{12bp}
\begin{multicols}{2}
\section*{1. 引言（原创示例）}
资料可能分布在正文、图表及文末独立框文中。完整呈现首先需要确定实际提供的内容，再核对其阅读顺序和边界\CVRCite{1}。

本页展示黑色通栏题名、粗体作者、浅米底结构化摘要及同页双栏开篇。普通刊名字样替代官方品牌标识，不使用源论文的全文、原图或个人下载信息。
\end{multicols}
\vfill\noindent\rule{\textwidth}{.4bp}\par
{\fontsize{7bp}{9bp}\selectfont *通信作者、原刊版权和许可说明按所给版本填写；作者单位、下载机构标记和期刊版权不是同一信息。\par
本文件为原创中性版式演示，不构成真实论文。仅供学术交流、出版物版权归原作者及出版商所有。\par}
\clearpage
\newgeometry{left=30bp,right=25.276bp,top=52bp,bottom=42bp,headheight=22bp,headsep=16.255bp,footskip=22bp}
\begin{multicols}{2}
\section*{2. 方法（结构示例）}
\CVRBody
\subsection*{2.1 记录与范围}
原刊主文页码与PDF物理页分别记录；收稿、接受、提前在线和正式刊期是不同字段。补充材料仅被引用时，不能视为已经取得或完成翻译。
\section*{3. 图文与表格对应}
本页通过原创矢量示意图及中性资料表展示通栏模块。实际论文的主图应靠近相关讨论，并在书目之前排完\CVRCite{2}。
\end{multicols}\vspace{5bp}
\noindent\begin{minipage}{\textwidth}
\begin{center}\begin{minipage}{.72\textwidth}\DemoVector\end{minipage}\end{center}
{\fontsize{7.6bp}{9.2bp}\selectfont\textbf{图1。原创记录核对示意（中性占位，无实验数据）。}
\textbf{A}，输入节点表示题名、摘要、正文与图表；\textbf{B}，处理节点表示范围、限定及单位的核对；\textbf{C}，输出节点表示正文、图注、书目及来源信息。箭头只表示本例的阅读顺序，不表示生物机制、因果关系或数量变化。图形由模板自行绘制，无真实数据、样本量、误差线或比例尺。实际原图保留英文时，应提供完整中文标签释义。\par}
\end{minipage}\vspace{12bp}
\begingroup\setlength{\fboxsep}{6bp}\noindent\colorbox{CVRCream}{\begin{minipage}{\dimexpr\textwidth-12bp\relax}
{\bfseries\color{CVRRed}表1}\quad\textbf{资料记录字段（原创中性占位）}\par\vspace{6bp}
{\fontsize{7.8bp}{10.2bp}\selectfont\arrayrulecolor{CVRRed}
\begin{tabularx}{\linewidth}{@{}>{\raggedright\arraybackslash}p{.21\linewidth}XXX@{}}\hline
\rowcolor{white}\textbf{记录} & \textbf{来源} & \textbf{对应内容} & \textbf{核对字段}\\\hline
\rowcolor{white}文字记录（占位） & 来源字段待填 & 正文与方法 & 范围、段落与限定\\
\rowcolor{white}图形记录（占位） & 来源字段待填 & 面板与图注 & 标签、单位与署名\\
\rowcolor{white}文末记录（占位） & 来源字段待填 & 声明与出处 & 身份、编号与链接\\\hline
\end{tabularx}\par\vspace{4bp}
表中各项均为原创排版占位，无真实材料、参数或研究数据。正式长表须逐行保留信息及全部脚注，必要时使用可续页结构。\par}
\end{minipage}}\endgroup\vspace{8bp}
\begin{multicols}{2}
\subsection*{3.1 通栏模块后的正文}
图注与表格应恢复完整阅读顺序，不能为了控制页数删去尾句。红色主节、黑色次节、蓝色数字上标及紧凑无衬线图注为本刊核验版式的特征。
\subsection*{3.2 证据与限定}
正文数据与排版样式分别维护；只改版式时复用已核对内容并检查指纹。图表裁区应逐幅测量，不继承其他文章的裁剪坐标。
\end{multicols}
\clearpage
\begin{multicols}{2}
\section*{4. 讨论（原创示例）}
\CVRBody
\subsection*{4.1 解释边界}
原文的相关、因果、假设及潜在用途具有不同强度，译文应保持这些区别。本例只演示版式，所有文字均不是实验结论或临床建议。
\subsection*{声明（原创占位）}
致谢、资助、作者贡献、利益冲突、数据与代码可用性及许可说明按来源分别核对。本示例没有真实作者、经费、数据或利益关系。
\subsection*{参考文献（原创占位）}
{\fontsize{7.6bp}{10bp}\selectfont\setlength{\parindent}{0pt}\interlinepenalty=10000
\hypertarget{cvr-ref:1}{}1. Author Placeholder A. Complete record inventory. Original neutral reference example; no publication or DOI assigned.\par\smallskip
\hypertarget{cvr-ref:2}{}2. Author Placeholder B. Correspondence between source and presentation. Original neutral reference example; no publication or DOI assigned.\par}
\end{multicols}\vspace{14bp}
\begingroup\setlength{\fboxsep}{8bp}\setlength{\fboxrule}{.5bp}\noindent\fcolorbox{black}{CVRGrey}{\begin{minipage}{\dimexpr\textwidth-17bp\relax}
{\bfseries\color{CVRRed}\fontsize{12.5bp}{15bp}\selectfont Translational perspective\quad 转化展望（原创占位）\par}\vspace{4bp}
本框展示位于参考文献之后的独立转化展望。正式译稿应完整翻译原作者对研究意义、后续验证和转化范围的阐述，并保留其中的不确定性与限定。不能用译者总结替代原框文，也不能因其位于书目之后而遗漏。本段不报告任何真实机制、效应或转化用途。灰底、黑框及红色框题根据所给版本核验；不将其他期刊的图形摘要结构套入本页。
\end{minipage}}\endgroup\vspace{14bp}
\section*{原文信息（占位字段）}
作者；英文原题；期刊；年份；卷期；页码或文章号；DOI；明确提供的提前在线日期和正式刊期。所有字段按所给版本逐项填写，此处没有对应的真实论文，不从文件名或DOI年份推填缺项。\par
本文件为Cardiovascular Research风格的非官方中性样板。普通刊名、字体回退及中文分页为有意适配，不代表该刊所有文体，也不承诺逐像素复刻。正文图位于书目前，独立转化展望保留在书目之后；最后以原文信息结束。\par
仅供学术交流、出版物版权归原作者及出版商所有。
\end{document}
